% Spanish -- One-page quick notes
% Compile: pdflatex spanish-only.tex
% Print this page only for Spanish exam

\documentclass[10pt,landscape,a4paper]{article}
\usepackage[utf8]{inputenc}
\usepackage[T1]{fontenc}
\usepackage{microtype}
\usepackage[left=1.2cm,right=1.2cm,top=1cm,bottom=1cm]{geometry}
\usepackage{multicol}
\usepackage{parskip}
\usepackage{xcolor}
\usepackage{enumitem}
\usepackage{booktabs}
\definecolor{accent}{RGB}{34,139,34}
\setlist{nosep,leftmargin=*,itemsep=0.2em}
\setlength{\parskip}{0.4em}
\setlength{\columnsep}{1em}
\raggedright

\begin{document}
\normalsize

\begin{center}
{\LARGE\bfseries\color{accent} Spanish -- Quick Notes}\\[0.3em]
\small One page. Bring only this for Spanish exam.
\end{center}
\vspace{0.5em}

\begin{multicols}{2}

\subsection*{\color{accent} If they ask... $\to$ Think...}
\begin{itemize}[leftmargin=*,nosep]
\item \textbf{Por vs para?} Para + inf (purpose). Porque + verb (reason). Por + noun (for/because of).
\item \textbf{C o qu?} qu before e, i. qui\'en = qu (not q). 6 = seis (not sies).
\item \textbf{Nacionalidad?} Agreement m/f. -o/-a. -ista, -ense invariable.
\item \textbf{Pronunciaci\'on?} Carlos [k], Garc\'ia [g], Javier [x], Celia [s].
\item \textbf{Listening?} nombre, origen, nacionalidad, edad, profesi\'on.
\end{itemize}

\subsection*{\color{accent} Preguntas (match question $\to$ answer)}
\begin{center}
\small
\begin{tabular}{@{}ll@{}}
\toprule
\textbf{Pregunta} & \textbf{Respuesta} \\
\midrule
¿C\'omo te llamas? & Francisco / Me llamo... \\
¿Cu\'antos a\~nos tienes? & Treinta y cinco. \\
¿Cu\'al es tu apellido? & Mart\'inez. \\
¿De d\'onde eres? & De Valencia. \\
¿En qu\'e trabajas? & Soy camarero. \\
\bottomrule
\end{tabular}
\end{center}

\subsection*{\color{accent} Por vs Para vs Porque}
\begin{center}
\small
\begin{tabular}{@{}ll@{}}
\toprule
\textbf{Use} & \textbf{Example} \\
\midrule
Para + infinitive & para hablar, para practicar \\
Porque + clause & porque quiero ver \\
Por + noun & por la comida, por el K-pop \\
\bottomrule
\end{tabular}
\end{center}

\subsection*{\color{accent} Ser \& Tener}
Ser: soy, eres, es, somos, sois, son. Tener: tengo, tienes, tiene, tenemos, ten\'eis, tienen. Llamarse: me llamo, te llamas, se llama... \textbf{Agreement} with person (m/f).

\subsection*{\color{accent} N\'umeros}
0--15: cero, uno... cinco, \textbf{seis}, siete... quince. 16--19: diecis\'eis, diecisiete... 20--50: veinte, treinta, cuarenta, cincuenta.

\subsection*{\color{accent} 5 lenguas Espa\~na}
Espa\~nol, euskera, catal\'an, aran\'es, gallego. Rom\'anicas: esp, cat, gal, aran\'es. Euskera: origen incierto.

\subsection*{\color{accent} Actividades}
Ver: pel\'iculas, series. Escuchar: m\'usica, radio. Leer: libros. Escribir: diario. Practicar: con nativos, pronunciaci\'on. Hacer: ejercicios.

\subsection*{\color{accent} Pronunciaci\'on}
Carlos [k]: c before a,o,u; Q. Garc\'ia [g]: G before a,o,u. Javier [x]: J; G before e,i. Celia [s]: c before e,i; Z.

\subsection*{\color{accent} Listening -- Una fiesta}
Ana: not Argentina, not cocinera. Elvira: colombiana. Stephanie: not Dutch. Jos\'e Antonio: profesor de espa\~nol.

\subsection*{\color{accent} Common mistakes}
6 = seis. qui\'en = qu. Gapfill: check gender (mi amiga $\to$ -a).

\end{multicols}
\end{document}
